\usepackage[numbers]{natbib} 

% Select input encodung, usually utf8 is the best choice, on windows, \usepackage[latin1]{inputenc} maybe required
\usepackage[utf8]{inputenc}
\usepackage[T1]{fontenc}
\usepackage[english]{babel}
\usepackage{csquotes}
\usepackage{xcolor}

\MakeOuterQuote{"} % Damit ist es möglich, " " zu verwenden ohne Umlaut zu erzeugen
\defaulthyphenchar=127 % Dadurch werden auch Wörter mit Bindestrich getrennt, die schon Bindestriche enthalten.

% Todo notes
\usepackage{todonotes}
\newcommand{\question}[2][]{%
    \todo[color=blue!40, #1]{#2}%
}

% geometry
\usepackage{framed}
\usepackage[bindingoffset=1cm, left=2.5cm, right=2.5cm, top=2.5cm, bottom=2.5cm]{geometry}

% Headline
\usepackage{fancyhdr}
\pagestyle{fancy}
\renewcommand{\chaptermark}[1]{\markboth{\thechapter\ #1}{}}
\lhead{\leftmark} \rhead{\thepage}
\cfoot{}
\fancypagestyle{plain}{}

\RedeclareSectionCommand[beforeskip=1.5cm,afterskip=1cm]{chapter}

% Colors
\usepackage{color}
\usepackage{colortbl}

% Tables
\usepackage{tabularx}
\usepackage{multirow}
\setlength{\tabcolsep}{4pt}

% Drawing graphs etc.
\usepackage{pgf}
\usepackage{tikz}
\usetikzlibrary{shapes.geometric,arrows,automata}

\tikzstyle{startstop} = [rectangle, rounded corners, 
minimum width=3cm, 
minimum height=1cm,
text centered, 
draw=black, 
fill=red!30]

\tikzstyle{io} = [trapezium, 
trapezium stretches=true, % A later addition
trapezium left angle=70, 
trapezium right angle=110, 
minimum width=3cm, 
minimum height=1cm, text centered, 
draw=black, fill=blue!30]

\tikzstyle{process} = [rectangle, 
minimum width=3cm, 
minimum height=1cm, 
text centered, 
text width=3cm, 
draw=black, 
fill=orange!30]

\tikzstyle{decision} = [diamond, 
minimum width=3cm, 
minimum height=1cm, 
text centered, 
draw=black, 
fill=green!30]
\tikzstyle{arrow} = [thick,->,>=stealth]

% Footnotes
\usepackage{footmisc}

\usepackage{xspace}
\newcommand{\sic}{[\acs{sic}]\xspace}

% math
\usepackage{amsmath}
\usepackage{amssymb}
\usepackage{siunitx}

% lists
\usepackage{paralist}

% Figures
\usepackage{graphicx, wrapfig}
\usepackage{subcaption}

% Hyperlinks
\usepackage[hyphens]{url}
\usepackage{hyperref}
\hypersetup{colorlinks, citecolor=black, linkcolor=black, urlcolor=black}

% Minted
\usepackage[chapter]{minted}
%\usemintedstyle{xcode}
\setminted{frame=single,tabsize=2,linenos,autogobble}

\newmintinline[code]{text}{breaklines}

\newminted[mdcodeblock]{md}{autogobble,frame=none,linenos=false,breaklines}

% list of abbreviations
\usepackage[printonlyused]{acronym}

% Set line pitch
\usepackage{setspace}
\onehalfspacing              % anderthalbzeilig (oder auch \doublespace)

%fancyBox
%\usepackage{fancybox}

% Layout corrections
\clubpenalty = 10000
% Layout corrections
\widowpenalty = 10000
\displaywidowpenalty = 10000

% Figures
\usepackage{caption}
\usepackage[hypcap=true,labelformat=simple]{subcaption}
\renewcommand{\thesubfigure}{(\alph{subfigure})}
\usepackage{placeins}

% Tables
\usepackage{booktabs}
\usepackage{enumitem}
\usepackage{multirow}
\usepackage{rotating} 

% Texts
\usepackage{seqsplit}

% Frequently used column types
\newcolumntype{C}[1]{>{\centering\arraybackslash}p{#1}} % centering column type with fixed width
\newcolumntype{R}[1]{>{\raggedleft\arraybackslash}p{#1}} % right aligned column type with fixed width
\newcolumntype{L}[1]{>{\raggedright\arraybackslash}p{#1}} % left aligned column type with fixed width

% Shortcuts for referencing floats:
\newcommand{\fig}[1]{\figurename~\ref{#1}} %shortcut for a figure reference
\newcommand{\tab}[1]{Table~\ref{#1}} %shortcut for a table reference
\newcommand{\eq}[1]{(\ref{#1})} %shortcut for an equation reference
\newcommand{\lst}[1]{Listing~\ref{#1}} %shortcut for a listing reference
\newcommand{\sect}[1]{Section~\ref{#1}} %shortcut for a Section reference

% Personal shortcuts
\usepackage{etoolbox}

\newenvironment{note}[2]{%
  \begin{leftbar}%
  #1%
  \ifstrempty{#2}{}{ \cite{#2}}%
  \begin{itemize}%
}{%
  \end{itemize}%
  \end{leftbar}%
}

\newenvironment{unfinished}{%
  \begin{leftbar}%
}{%
  \end{leftbar}%
}

\usepackage{chngcntr}
\counterwithout{equation}{section}

\usepackage{xspace}
\newcommand{\benchmark}{\textit{Power Benchmark}\xspace}

\newcommand{\thesisrepo}{\url{https://github.com/Giulcoo/power-benchmark-thesis}}
\newcommand{\coderepo}{\url{https://github.com/Giulcoo/power-benchmark}}
\newcommand{\documentation}{https://github.com/Giulcoo/power-benchmark/wiki}

\newcounter{sublisting}[listing]
\renewcommand{\thesublisting}{\thelisting\alph{sublisting}}


% For creating overview tables for scores in experiment result
% #1 caption
% #2 label
% #3 table body (rows, already separated by & and \\, incl. \midrule)
\newcommand{\agentresulttable}[3]{%
  \begin{table}[htbp]
    \centering
    \caption{#1}
    \label{#2}
    \begin{tabular}{lrrr}
      \toprule
      Metric & Do-Nothing & Greedy & PPO \\
      \midrule
      #3
      \bottomrule
    \end{tabular}
  \end{table}
}